% block1_ftcs.tex
%
% Phase 3, Block 1: the explicit FTCS scheme for the Black-Scholes PDE.
%
% Self-contained writeup. Compiles standalone with pdflatex.

\documentclass[11pt,a4paper]{article}

\usepackage[T1]{fontenc}
\usepackage[utf8]{inputenc}
\usepackage{amsmath, amssymb, amsthm}
\usepackage{geometry}
\usepackage{hyperref}
\usepackage{enumitem}

\geometry{margin=1in}

\newtheorem{proposition}{Proposition}[section]
\newtheorem{theorem}[proposition]{Theorem}
\newtheorem{definition}[proposition]{Definition}
\newtheorem{remark}[proposition]{Remark}

\title{Phase 3, Block 1:\\The explicit FTCS scheme\\for the Black--Scholes PDE}
\author{Quant Finance Portfolio}
\date{}

\begin{document}
\maketitle

\section{Setting}

This block implements the Forward-Time, Centred-Space (FTCS) finite-
difference pricer for European calls and puts. We use the
constant-coefficient transformed Black--Scholes equation derived in
Block~0,
\begin{equation}
\frac{\partial V}{\partial \tau}
\;=\; \frac{\sigma^2}{2}\frac{\partial^2 V}{\partial x^2}
   + \mu\,\frac{\partial V}{\partial x}
   - r\, V,
\qquad \mu = r - \tfrac{\sigma^2}{2},
\label{eq:transformed}
\end{equation}
with $x = \ln(S/K)$, $\tau = T - t$, on the truncated domain
$[x_{\min}, x_{\max}] \times [0, T]$ with Dirichlet boundary
conditions specified in Block~0. The deliverable is a pricer,
validated against the Black--Scholes closed form and analysed in
terms of stability and convergence.

% =====================================================================
\section{The FTCS discretisation}
% =====================================================================

Let $V_j^n \approx V(x_j, \tau_n)$ on the uniform grid $x_j = x_{\min}
+ j\,\Delta x,\;\tau_n = n\,\Delta\tau$. Substituting the centred space
stencils and the forward-time stencil into \eqref{eq:transformed}:
\[
\frac{V_j^{n+1} - V_j^n}{\Delta\tau}
\;=\; \frac{\sigma^2}{2}\,\frac{V_{j+1}^n - 2\,V_j^n + V_{j-1}^n}{\Delta x^2}
+ \mu\,\frac{V_{j+1}^n - V_{j-1}^n}{2\,\Delta x}
- r\,V_j^n.
\]
All spatial differences are evaluated at the \emph{old} time level $n$.
This is what makes the scheme explicit: the new-time value $V_j^{n+1}$
is given directly as a linear combination of three old-time values,
without solving any system.

Multiplying by $\Delta\tau$ and isolating $V_j^{n+1}$ on the left,
\begin{equation}
\boxed{\;\;
V_j^{n+1} \;=\; a_-\,V_{j-1}^n \,+\, a_0\,V_j^n \,+\, a_+\,V_{j+1}^n
\;\;}
\label{eq:ftcs}
\end{equation}
with constant coefficients
\begin{equation}
a_-  = \alpha - \frac{\mu\,\Delta\tau}{2\,\Delta x},\quad
a_0  = 1 - 2\alpha - r\,\Delta\tau,\quad
a_+  = \alpha + \frac{\mu\,\Delta\tau}{2\,\Delta x},
\label{eq:coeffs}
\end{equation}
where $\alpha = \frac{\sigma^2}{2}\,\Delta\tau / \Delta x^2$ is the
Fourier number from Block~0. The coefficients do not depend on
$j$ or $n$, by virtue of the constant-coefficient structure of
\eqref{eq:transformed}: this is the structural pay-off of the
logarithmic transformation.

% =====================================================================
\section{Probabilistic interpretation}
% =====================================================================

The three coefficients sum to
\[
a_- + a_0 + a_+ \;=\; 1 - r\,\Delta\tau.
\]
This identity has a clean reading: if $a_-$, $a_0$, $a_+$ are
non-negative, define
\[
p_-  = \frac{a_-}{1 - r\Delta\tau},\quad
p_0  = \frac{a_0}{1 - r\Delta\tau},\quad
p_+  = \frac{a_+}{1 - r\Delta\tau},
\]
so that $p_-, p_0, p_+ \ge 0$ and $p_- + p_0 + p_+ = 1$. Then
\eqref{eq:ftcs} is exactly
\begin{equation}
V_j^{n+1} \;=\; (1 - r\,\Delta\tau)\,
\Bigl(p_-\,V_{j-1}^n + p_0\,V_j^n + p_+\,V_{j+1}^n\Bigr)
\;\approx\; e^{-r\Delta\tau}\,
\mathbb{E}^{\mathbb{Q}}\!\Bigl[\,V(x_j + \xi\,\Delta x,\,\tau_n)\,\Bigr],
\label{eq:prob-reading}
\end{equation}
where $\xi \in \{-1, 0, +1\}$ with $\mathbb{Q}$-probabilities
$(p_-, p_0, p_+)$. The factor $1 - r\,\Delta\tau$ is the
forward-Euler approximation to the discount factor $e^{-r\Delta\tau}$.

Equation \eqref{eq:prob-reading} says that FTCS is the discretisation
of the risk-neutral pricing identity
\[
V(x_j, \tau_{n+1})
\;=\; e^{-r\Delta\tau}\,\mathbb{E}^{\mathbb{Q}}\!\bigl[V(X_{n+1}, \tau_n)\,
\big|\, X_n = x_j\bigr]
\]
on a discrete state space, with $X_{n+1} - X_n$ taking the values
$\pm\Delta x, 0$. This is exactly the structure of a trinomial tree.
We will revisit this connection in Block~5: the Kamrad--Ritchken tree
is, after a precise reparametrisation, FTCS in disguise.

The reading also highlights why the non-negativity of the coefficients
matters: if any $p_\xi < 0$, the operation loses its interpretation as
a discrete expectation, and the scheme can amplify oscillations near
the strike kink in the payoff.

% =====================================================================
\section{Non-negativity of the coefficients}\label{sec:positivity}
% =====================================================================

Three conditions, one per coefficient.

\paragraph{Central coefficient.}
$a_0 \ge 0$ requires
\begin{equation}
\alpha \;\le\; \frac{1 - r\,\Delta\tau}{2}.
\label{eq:cfl-strict}
\end{equation}
For $r\,\Delta\tau \ll 1$ (always true in this project: $r \le 0.1$
and $\Delta\tau \le 10^{-2}$), \eqref{eq:cfl-strict} is essentially
the CFL condition $\alpha \le 1/2$ from Block~0.

\paragraph{Lateral coefficients.}
$a_\pm \ge 0$ requires $\bigl|\,\mu\,\Delta\tau/(2\Delta x)\bigr|\le
\alpha$, equivalently
\begin{equation}
\frac{|\mu|\,\Delta x}{\sigma^2} \;\le\; 1.
\label{eq:peclet}
\end{equation}
The dimensionless ratio $\mathrm{Pe}_h := |\mu|\Delta x / \sigma^2$ is
the \emph{grid Péclet number}: it measures the relative weight of
convection to diffusion at the scale of one grid cell. For the
parameter range of this project ($\sigma \in [0.1, 0.4]$,
$|\mu| \le 0.1$, $\Delta x \le 0.1$), $\mathrm{Pe}_h \le 1$ is
satisfied with comfortable margin.

\begin{remark}[Convection-dominated regimes]
When $\mathrm{Pe}_h > 1$ (extreme dividends, very small volatility,
very coarse grid), centred differences for the convective term
produce one negative lateral coefficient, the probabilistic reading
breaks, and the scheme exhibits localised oscillations near the
strike kink. The standard remedy is an \emph{upwind} discretisation
of the convective term: replace
$(V_{j+1} - V_{j-1}) / (2\Delta x)$ by $(V_j - V_{j-1})/\Delta x$ if
$\mu > 0$ or $(V_{j+1} - V_j)/\Delta x$ if $\mu < 0$. This restores
$L^\infty$ stability at the price of one order of accuracy in space
($O(\Delta x)$ instead of $O(\Delta x^2)$). Out of scope for Phase~3.
\end{remark}

\paragraph{Summary.}
The probability-positive (monotone) regime is
\begin{equation}
\boxed{\;\;\alpha \le \tfrac{1}{2}\;\;\text{and}\;\;\mathrm{Pe}_h \le 1.\;\;}
\label{eq:ppos}
\end{equation}
In this regime the scheme is monotone, and monotonicity implies
stability in the uniform $L^\infty$ norm:
\[
\max_j |V_j^{n+1}|
\;\le\; (1 - r\Delta\tau)\,\max_j |V_j^n|
\;\le\; \max_j |V_j^n|.
\]
This is a strictly stronger result than the $L^2$ stability of the
next section, and it is the practical reason the prices computed by
FTCS \emph{look smooth}: oscillations cannot grow.

% =====================================================================
\section{Von Neumann stability analysis}\label{sec:vN}
% =====================================================================

The von Neumann analysis of \eqref{eq:ftcs} for the full
convection-diffusion-reaction equation supplements the previous
section with a tight necessary-and-sufficient condition for $L^2$
stability.

Substituting the Fourier ansatz $V_j^n = g^n e^{ij\theta}$ into
\eqref{eq:ftcs} and using $e^{i\theta} + e^{-i\theta} = 2 - 4\sin^2(\theta/2)$,
$e^{i\theta} - e^{-i\theta} = 2 i\sin\theta$, the amplification factor is
\begin{equation}
g(\theta) \;=\; (1 - r\,\Delta\tau)
- 4\,\alpha\,\sin^2\!\tfrac{\theta}{2}
+ i\,\beta\,\sin\theta,
\qquad
\beta \;:=\; \frac{\mu\,\Delta\tau}{\Delta x}.
\label{eq:vn-factor}
\end{equation}
($\beta$ is the unsymmetrised Courant number; it differs from the
$\nu$ of Block~0 by a factor of $1/2$.) The factor is complex: the
real part is the diffusive contribution shifted by the reactive term;
the imaginary part is the convective contribution.

\begin{proposition}[Stability of FTCS]
Set $r = 0$ for clarity (the reactive term shifts the analysis
trivially and never contributes to instability when $r \ge 0$ and
$\Delta\tau \le 1/r$). The condition $|g(\theta)|^2 \le 1$ for all
$\theta \in [-\pi, \pi]$ is equivalent to
\begin{equation}
\alpha \;\le\; \tfrac{1}{2}
\qquad\text{and}\qquad
\beta^2 \;\le\; 2\,\alpha.
\label{eq:vn-conds}
\end{equation}
\end{proposition}

\begin{proof}
With $s := \sin^2(\theta/2) \in [0, 1]$, $\sin^2\theta = 4 s (1-s)$.
Then
\[
|g|^2 = (1 - 4\alpha s)^2 + 4\beta^2 s(1-s),
\]
and after expansion,
\[
|g|^2 - 1 = s\,\bigl[(16\alpha^2 - 4\beta^2)\,s + (4\beta^2 - 8\alpha)\bigr].
\]
The bracket is linear in $s$. The condition $|g|^2 \le 1$ for all
$s \in (0, 1]$ is equivalent to the bracket being non-positive at
both endpoints $s = 0$ and $s = 1$ (linear function on a compact
interval attains its extrema at the endpoints). At $s = 0$:
$4\beta^2 - 8\alpha \le 0 \Leftrightarrow \beta^2 \le 2\alpha$. At
$s = 1$: $16\alpha^2 - 8\alpha \le 0 \Leftrightarrow \alpha \le 1/2$.
\qed
\end{proof}

\paragraph{The two conditions are not equivalent to monotonicity.} The
von Neumann conditions \eqref{eq:vn-conds} are about $L^2$
stability; the conditions \eqref{eq:ppos} are about monotonicity. The
relationship between them is:

\begin{itemize}[leftmargin=2em]
\item Both conditions share $\alpha \le 1/2$.
\item The von Neumann condition $\beta^2 \le 2\alpha$ rewrites as
$\mu^2 \Delta\tau \le \sigma^2$, a constraint on $\Delta\tau$. For
typical parameters ($\sigma = 0.2,\;\mu \le 0.1$): $\Delta\tau \le 4
\sigma^2/\mu^2 = 4$ years. Always satisfied.
\item The probability-positive condition $\mathrm{Pe}_h \le 1$
rewrites as $|\mu| \Delta x \le \sigma^2$, a constraint on $\Delta x$.
\end{itemize}
The two are different conditions on different variables. The
probability-positive constraint is strictly stronger: it implies $L^2$
stability via monotonicity, but the converse fails. In our parameter
range both conditions are satisfied; in convection-dominated regimes,
they decouple, and the von Neumann condition can hold without
monotonicity (giving stable but oscillatory schemes).

\paragraph{What dominates in practice.} The constraint that bites is
$\alpha \le 1/2$. The Péclet condition $\mathrm{Pe}_h \le 1$ is a
constraint on $\Delta x$ that is comfortably met for the project's
parameters. The reactive contribution $r\Delta\tau$ damps every Fourier
mode uniformly and never destabilises.

% =====================================================================
\section{The time-marching algorithm}
% =====================================================================

In pseudocode, the inner loop is
\begin{verbatim}
    V[j] for j = 0..N  <-  initial-condition vector at tau = 0
    for n = 0, 1, ..., M-1:
        V_new[0] = bc_lower[n+1]                          # Dirichlet
        V_new[N] = bc_upper[n+1]                          # Dirichlet
        for j = 1, ..., N-1:                              # interior nodes
            V_new[j] = a_minus * V[j-1]
                     + a_zero  * V[j]
                     + a_plus  * V[j+1]
        swap(V, V_new)
    return V                                              # V at tau = T
\end{verbatim}

Three implementation points worth explicit comment.

\paragraph{Double buffer.} The update of node $j$ at level $n+1$ uses
$V_{j-1}^n$ and $V_{j+1}^n$. If the implementation overwrites
$V_j$ in place, the value of $V_j$ used by the update of $V_{j+1}$
will be the new one rather than the old, breaking the scheme. Two
buffers $V$ and $V_{\text{new}}$ are needed, swapped after each step.

\paragraph{Boundary conditions are not stencil applications.} The
boundary nodes $j = 0$ and $j = N$ are fixed by the Dirichlet
conditions of Block~0; the stencil is applied only to interior
nodes $j = 1, \dots, N-1$. Trying to apply the stencil at $j = 0$
would require a phantom node $j = -1$ outside the domain.

\paragraph{Recovery of the price at $S_0$.} The grid is in $x =
\ln(S/K)$. The pricer must return $V(S_0, 0)$ at the user's spot
$S_0$, which is generally not on a grid node. The $V$ array at
$\tau = T$ contains $V_j^M$ for $j = 0, \dots, N$, and the price is
recovered by interpolation at $x_0 = \ln(S_0/K)$. Linear
interpolation is enough: it adds an $O(\Delta x^2)$ error consistent
with the spatial discretisation order.

% =====================================================================
\section{Convergence and computational cost}
% =====================================================================

\paragraph{Local truncation error.} By Taylor expansion of every
stencil,
\[
\tau_{j}^n \;=\; O(\Delta\tau) + O(\Delta x^2).
\]
First order in $\Delta\tau$ from the forward-Euler step; second order
in $\Delta x$ from the centred space stencils.

\paragraph{Global error and convergence rate.} By the Lax equivalence
theorem (Block 0, §7), consistency plus stability gives convergence,
\[
\bigl\|V_h - V\bigr\|_\infty \;=\; O(\Delta\tau) + O(\Delta x^2).
\]
Under the CFL constraint $\Delta\tau \le \Delta x^2 / \sigma^2$, both
contributions are $O(\Delta x^2)$, so the scheme converges
quadratically in $\Delta x$. There is no parameter regime in which
FTCS achieves better than $O(\Delta x^2)$ convergence: the temporal
order is locked at $1$, and the spatial order is locked at $2$, but
the CFL constraint balances the two automatically.

\paragraph{Computational cost.} Each time step does $O(N)$ work
(one stencil application per interior node). The number of time
steps is, by CFL,
\[
M \;\ge\; \frac{T\,\sigma^2}{\Delta x^2} \;\sim\; N^2.
\]
Total work: $O(N \cdot M) = O(N^3)$. Doubling $N$ multiplies the
cost by $8$. This $N^3$ scaling is the price tag of explicit time
stepping; Blocks~2 (BTCS) and~3 (Crank--Nicolson) will sidestep it
by allowing $\Delta\tau \sim \Delta x$ at the cost of solving a
tridiagonal system per step.

% =====================================================================
\section{Implementation outline}
% =====================================================================

The Python module \texttt{quantlib.ftcs} exposes
\texttt{ftcs\_european\_call} and \texttt{ftcs\_european\_put}, both
with signature \texttt{(S, K, r, sigma, T, *, N, M, n\_sigma=4.0)}.
The C++ counterpart in the \texttt{quant::pde} namespace has the
same arguments. Both validate that $\alpha \le 1/2$ at construction
and raise an error otherwise (Python: \texttt{ValueError}; C++:
\texttt{std::invalid\_argument}). The validation is strict and not
configurable: the pricer will not silently produce nonsense from a
violated CFL.

The validation tests (Python: \texttt{validate\_ftcs.py}; C++:
\texttt{test\_ftcs.cpp}) check four properties:
\begin{enumerate}[leftmargin=2em]
\item agreement with the Black--Scholes closed form on a grid of
parameters, with tolerance scaled to $\Delta x^2$;
\item the empirical convergence rate, by running on $(N, M)$ and
$(2N, 4M)$ and checking that the error ratio is close to $4$;
\item put-call parity: $C - P = S - K e^{-rT}$ at every spot, to
within twice the discretisation tolerance;
\item that violating CFL ($\alpha > 1/2$) raises the appropriate
exception.
\end{enumerate}

A separate diagnostic script, \texttt{demo\_ftcs\_explosion.py},
disables the CFL check and runs the scheme at $\alpha = 0.6$, to
visualise the sawtooth instability mode. The output is documented in
the manual personal as the Phase~3 numerical lesson.

% =====================================================================
\section*{References}
% =====================================================================

\begin{thebibliography}{9}

\bibitem{TavellaRandall2000}
D. Tavella and C. Randall.
\emph{Pricing Financial Instruments: The Finite Difference Method.}
Wiley, 2000. \S3.1--3.3 (FTCS for BS, the CFL constraint specific to
the equation), \S4.2 (convergence).

\bibitem{Seydel2017}
R. Seydel.
\emph{Tools for Computational Finance.}
6th ed., Springer, 2017. \S4.2.1 (FTCS), \S4.2.4 (stability analysis).

\bibitem{Strikwerda2004}
J.~C. Strikwerda.
\emph{Finite Difference Schemes and Partial Differential Equations.}
2nd ed., SIAM, 2004. \S3.4--3.5 (von Neumann analysis for the full
CDR equation).

\bibitem{MortonMayers2005}
K.~W. Morton and D.~F. Mayers.
\emph{Numerical Solution of Partial Differential Equations.}
2nd ed., CUP, 2005. \S2.7--2.10 (Lax equivalence and consistency
analysis).

\bibitem{Hull2017}
J. Hull.
\emph{Options, Futures, and Other Derivatives.}
10th ed., Pearson, 2017. \S21.4 (FTCS \(\leftrightarrow\) trinomial
tree equivalence, from the tree side).

\bibitem{KamradRitchken1991}
B. Kamrad and P. Ritchken.
\emph{Multinomial approximating models for options with $k$ state
variables.}
Management Science 37(12), 1640--1652, 1991. The original trinomial
tree paper; will be revisited in Block~5.

\end{thebibliography}

\end{document}
